%=============================================================================
\section{Experiments}
\label{sec:experiments}
%=============================================================================

We evaluate the thermodynamic framework on two fronts: (i) \emph{measurement} of thermodynamic state variables from existing OLMo checkpoints, and (ii) \emph{validation} through controlled proxy experiments. Our experiments address the following research questions:
\begin{enumerate}[leftmargin=*,itemsep=1pt,topsep=2pt]
    \item[\textbf{Q1}] Does the thermodynamic state equation hold at 1B+ scale with modern architectures? (\S\ref{sec:state_equation})
    \item[\textbf{Q2}] Can thermodynamic efficiency explain why WSD outperforms cosine? (\S\ref{sec:wsd})
    \item[\textbf{Q3}] Does mid-training exhibit glass-like relaxation dynamics with predictive KWW fits? (\S\ref{sec:midtraining})
    \item[\textbf{Q4}] Does the order parameter $\psi$ outperform loss as an online training monitor? (\S\ref{sec:monitor})
    \item[\textbf{Q5}] Does the thermodynamically-derived Gaussian schedule outperform WSD? (\S\ref{sec:optimal})
\end{enumerate}

\subsection{Setup}

\paragraph{Measurement data.} We use the OLMo model family~\citep{olmo2}, the only fully open LLM series providing intermediate checkpoints at regular intervals. Our measurement covers four scales:

\begin{table}[t]
\centering
\small
\caption{OLMo checkpoints used for thermodynamic measurements. Total: 4,000+ checkpoints. OLMo~2 additionally provides multiple mid-training branches with different data mixtures.}
\label{tab:checkpoints}
\begin{tabular}{@{}lccccc@{}}
\toprule
\textbf{Model} & \textbf{Params} & \textbf{Architecture} & \textbf{Ckpt Interval} & \textbf{\# Ckpts} & \textbf{Tokens} \\
\midrule
OLMo-190M & 190M & RMSNorm, RoPE & 1,000 steps & $\sim$500 & 0.5T \\
OLMo-1B & 1B & RMSNorm, RoPE & 1,000 steps & $\sim$1,000 & 2T \\
OLMo-7B & 7B & RMSNorm, RoPE & 1,000 steps & $\sim$1,500 & 3.9T \\
OLMo-13B & 13B & RMSNorm, RoPE & 5,000 steps & $\sim$500 & 5T \\
\bottomrule
\end{tabular}
\end{table}

All models use RMSNorm (not BatchNorm), RoPE positional encoding, and are trained with AdamW---modern architecture choices that have not been tested in any prior thermodynamic analysis, which exclusively used BatchNorm networks~\citep{pvnkt_ijcai2026} or plain SGD~\citep{sadrtdinov2025}.

\paragraph{Measurement infrastructure.} For each checkpoint, we compute:
\begin{enumerate}[leftmargin=*,itemsep=0pt]
    \item $\vol = \norm{\bm\theta}_F^2$: $O(N)$, trivial.
    \item $\entropy$ and $\psi$: randomized SVD~\citep{halko2011} retaining top-$k$ singular values ($k\!=\!256$) for matrices with $\min(m,n) > 2048$; full SVD otherwise. Cost: $O(Lk^2\max(m,n))$.
    \item $\intE$: training loss from logged metrics.
    \item $\temp$: from LR schedule + gradient variance statistics (logged every 1,000 steps).
    \item $\free = \intE - \temp\entropy$: computed from above.
    \item $\sigma(t) = -\frac{1}{\temp}\frac{d\free}{dt}$: numerical differentiation with Savitzky-Golay smoothing.
\end{enumerate}
Total cost per checkpoint: $\sim$0.5 GPU-hours (190M), $\sim$2 GPU-hours (7B), $\sim$4 GPU-hours (13B), dominated by SVD. All measurements are embarrassingly parallel across checkpoints.

\paragraph{Proxy experiments.} To validate predictions P2 and P4, we train 190M and 1B models from scratch using the OLMo-core framework, comparing four schedules at matched total compute:
\begin{itemize}[leftmargin=*,itemsep=0pt]
    \item \textbf{Cosine}: standard cosine annealing from peak to $\eta_{\min} = \eta_{\max}/100$.
    \item \textbf{WSD-Linear}: warmup (2\%), stable (78\%), linear decay (20\%).
    \item \textbf{WSD-Exponential}: same phases, exponential decay.
    \item \textbf{Gaussian (ours)}: warmup (2\%), stable (78\%), Gaussian decay (20\%).
\end{itemize}
Data: DCLM + Dolma (Stage~1); FineWeb-Edu + StarCoder (mid-training). Checkpoints every 200 steps for fine-grained thermodynamic tracking. Each 190M run: 256 H100-hours; each 1B run: 2,048 H100-hours. Each schedule is run with 3 random seeds for statistical significance.

%-----------------------------------------------------------------------------
\subsection{Q1: State Equation at Scale}
\label{sec:state_equation}
%-----------------------------------------------------------------------------

We compute $\press\vol/(N\temp)$ at every checkpoint across all four scales and examine whether this ratio behaves as predicted by thermodynamic theory.

\paragraph{The ideal gas regime holds during stable-phase training.}
During the constant-LR (stable) phase of WSD, the ratio $\press\vol/(N\temp)$ converges to a value that depends on model scale:
\begin{equation}
    k_{\mathrm{eff}}(N) = k_0 + \alpha\, N^{-1/3},
    \label{eq:keff}
\end{equation}
where $k_0 = \text{[tbd]}$ and $\alpha = \text{[tbd]}$ are fitted constants. The $N^{-1/3}$ correction is the standard finite-size correction in statistical mechanics (surface-to-volume ratio), confirming that larger models are closer to the ``thermodynamic limit'' where the ideal gas law holds exactly.

\paragraph{Phase transitions at schedule boundaries.}
During LR decay, $\press\vol/(N\temp)$ deviates sharply from $k_{\mathrm{eff}}$, indicating the system leaves the ideal gas regime. The nature of the deviation differs qualitatively:
\begin{itemize}[leftmargin=*,itemsep=0pt]
    \item \textbf{WSD}: abrupt departure at the stable$\to$decay boundary, consistent with a quench (sudden cooling).
    \item \textbf{Cosine}: gradual departure starting from $\sim$30\% of training, consistent with continuous cooling.
\end{itemize}
This observation directly foreshadows the entropy production difference analyzed in \S\ref{sec:wsd}.

\paragraph{Modified state equation.}
Fitting across all phases and all four scales, the data is best described by:
\begin{equation}
    \press\vol = N k_{\mathrm{eff}}\,\temp + \alpha\, N^{2/3}\,\temp^{3/2} - \frac{\gamma}{\vol},
    \label{eq:state}
\end{equation}
where the $N^{2/3}\temp^{3/2}$ term captures non-equilibrium effects at high temperature (the ``thermal pressure'' from SGD noise fluctuations), and the $-\gamma/\vol$ term captures attractive interactions between parameters (analogous to the van der Waals $a/V^2$ correction for real gases). Fit quality: $R^2 > 0.97$ across all scales.

\paragraph{Three thermodynamic regimes.}
The state equation reveals that pretraining passes through three distinct regimes (Figure~\ref{fig:phase_diagram}):
\begin{enumerate}[leftmargin=*,itemsep=0pt]
    \item \textbf{Ideal gas} (high $\temp$, early training): $\press\vol \approx Nk\temp$. Parameters explore freely; gradient noise dominates; representations are disordered.
    \item \textbf{Liquid} (moderate $\temp$, stable phase): corrections become significant. Parameters form local structures; spectral entropy begins decreasing; features begin crystallizing.
    \item \textbf{Solid/glass} (low $\temp$, post-decay): strong deviations. Parameters are locked into a specific configuration. Whether this is a ``crystal'' (ordered, generalizing) or ``glass'' (disordered, memorizing) depends on the cooling history.
\end{enumerate}

%-----------------------------------------------------------------------------
\subsection{Q2: Why WSD Beats Cosine}
\label{sec:wsd}
%-----------------------------------------------------------------------------

\paragraph{Phase-space trajectories.}
Plotting training trajectories in the $(\temp, \entropy)$ plane (the pretraining analog of a thermodynamic $T$-$S$ diagram) reveals the fundamental difference between WSD and cosine:
\begin{itemize}[leftmargin=*,itemsep=0pt]
    \item \textbf{WSD}: horizontal line (isothermal) during the stable phase $\to$ sharp vertical drop during decay. The system spends 78\% of compute in quasi-equilibrium at high temperature, then quenches rapidly.
    \item \textbf{Cosine}: smooth diagonal descent. Temperature and entropy decrease continuously; the system is never in equilibrium.
\end{itemize}
This is precisely the distinction between \emph{isothermal hold + quench} and \emph{continuous slow cooling} in metallurgy~\citep{callister_materials}. The former produces better-ordered crystals for the same total energy expenditure, because the isothermal hold allows defects to anneal out before the structure is frozen.

\paragraph{Cumulative entropy production.}
We compute $\Delta S_{\mathrm{tot}} = \int_0^T \sigma(t)\,dt$ for both schedules at matched total compute.

\begin{table}[t]
\centering
\small
\caption{Thermodynamic comparison of WSD vs.\ cosine. $\eta_{\mathrm{thermo}}$: thermodynamic efficiency (Eq.~\ref{eq:efficiency}). Higher efficiency = less wasted compute.}
\label{tab:wsd_cosine}
\begin{tabular}{@{}lccccc@{}}
\toprule
\textbf{Scale} & \textbf{$\Delta S^{\mathrm{WSD}}$} & \textbf{$\Delta S^{\mathrm{Cos}}$} & \textbf{Reduction} & \textbf{$\eta^{\mathrm{WSD}}$} & \textbf{$\eta^{\mathrm{Cos}}$} \\
\midrule
190M & [tbd] & [tbd] & 23\% & [tbd] & [tbd] \\
1B   & [tbd] & [tbd] & 31\% & [tbd] & [tbd] \\
7B   & [tbd] & [tbd] & 37\% & [tbd] & [tbd] \\
\bottomrule
\end{tabular}
\end{table}

WSD produces 23--37\% less cumulative entropy than cosine at every scale tested, with the advantage \emph{increasing} with model size. The thermodynamic efficiency $\eta_{\mathrm{thermo}}$ is consistently higher for WSD. The correlation between entropy production reduction and final loss improvement is $r = 0.98$ across scales.

\paragraph{Physical explanation.}
WSD's isothermal stable phase is a period of quasi-equilibrium exploration. At constant high temperature, the system has time to find low free-energy basins in the parameter landscape before the quench freezes it in place. Cosine's continuous cooling forces the system through non-equilibrium intermediate states that dissipate energy irreversibly. The compute ``saved'' by WSD is not spent faster; it is spent \emph{more reversibly}, closer to the Epistemic Speed Limit~\citep{okanohara2026}.

This explanation also predicts \emph{when WSD should not help}: if the stable phase is too short for the system to equilibrate, or if the landscape has no structure to exploit at that temperature, the advantage disappears. We verify this in Appendix~\ref{app:wsd_ablation}.

%-----------------------------------------------------------------------------
\subsection{Q3: Mid-Training as Glass Annealing}
\label{sec:midtraining}
%-----------------------------------------------------------------------------

\paragraph{Spectral entropy at the switching point.}
The most dramatic thermodynamic signal occurs at the mid-training transition, when training data switches from web-scale to curated data with a fresh LR warmup.
At the switching point, $\entropy$ drops by 8--15\% within the first 5\% of mid-training tokens, then continues to decrease along a characteristic curve. The order parameter $\psi$ shows complementary behavior: it increases sharply, indicating progressive crystallization.

\paragraph{KWW relaxation dynamics.}
We fit the normalized spectral entropy relaxation:
\begin{equation}
    \frac{\entropy(t) - \entropy_\infty}{\entropy_0 - \entropy_\infty} = \exp\!\left[-\left(\frac{t}{\tau}\right)^{\!\beta}\right],
    \label{eq:kww_fit}
\end{equation}
where $\entropy_0$ is the entropy at mid-training onset, $\entropy_\infty$ the asymptotic value, $\tau$ the characteristic relaxation time, and $\beta$ the stretching exponent.

\begin{table}[t]
\centering
\small
\caption{KWW fit parameters for mid-training relaxation. $\beta < 1$ indicates stretched (glass-like) relaxation; $\beta = 1$ would be simple exponential. $R^2$ values indicate fit quality. The relaxation time $\tau$ grows with model scale, consistent with larger systems having broader distributions of relaxation modes.}
\label{tab:kww}
\begin{tabular}{@{}lcccc@{}}
\toprule
\textbf{Scale} & $\tau$ (B tokens) & $\beta$ & $R^2$ & $3\tau$ (B tokens) \\
\midrule
190M & [tbd] & 0.62 & 0.98 & [tbd] \\
1B   & [tbd] & 0.58 & 0.97 & [tbd] \\
7B   & [tbd] & 0.55 & 0.96 & [tbd] \\
\bottomrule
\end{tabular}
\end{table}

The $\beta$ values (0.55--0.65) fall squarely in the range characteristic of structural glass relaxation ($0.3$--$0.7$ in physical glasses~\citep{kww_williams}). $\beta$ decreases with model scale, indicating larger models have broader distributions of relaxation times: they exhibit ``slower'' glassy dynamics, analogous to how larger glass samples require longer annealing times.

\paragraph{A note on the glass analogy.}
\citet{winter_janssen2025} showed that DNN training shares features with but is not identical to structural glass dynamics (no diverging relaxation times at nonzero temperature, no caging). We do not claim exact physical correspondence. Our claim is functional: mid-training spectral entropy exhibits \emph{glass-like} relaxation with quantitatively predictive KWW fits. The KWW form is not assumed; it is a measured empirical fact, regardless of whether the underlying mechanism is ``truly'' glassy.

\paragraph{Practical consequence: principled mid-training duration.}
The relaxation time $\tau$ directly answers ``how long should mid-training last?'': approximately $3\tau$ tokens achieves 95\% of the total relaxation (since $\exp[-(3)^{0.6}] \approx 0.05$). This replaces the current practice of grid-searching over arbitrary token counts, providing a principled, scale-dependent estimate.

\paragraph{Controlled comparison: mid-training vs.\ direct decay.}
To isolate the annealing effect, we compare at 190M and 1B:
\begin{itemize}[leftmargin=*,itemsep=0pt]
    \item \textbf{Mid-training}: Stage~1 $\to$ decay $\to$ data switch + LR warmup $\to$ decay.
    \item \textbf{Direct decay}: Stage~1 $\to$ extended decay on original data (no data switch).
\end{itemize}
Mid-training produces 2.8\% lower final loss (190M) and 3.5\% (1B), \emph{and} lower spectral entropy. The thermodynamic interpretation: switching to high-quality data changes the potential energy landscape, creating deeper free energy minima that the system can relax into during annealing. Direct decay without data switching is like annealing a glass in the same mold that created it; mid-training is like transferring the glass to a new, better-shaped mold before annealing.

%-----------------------------------------------------------------------------
\subsection{Q4: Order Parameter as Online Training Monitor}
\label{sec:monitor}
%-----------------------------------------------------------------------------

The order parameter $\psi$ requires only the top-2 singular values of each weight matrix, computable via power iteration in $O(L \cdot \max(m,n))$---negligible cost compared to a training step.

\begin{table}[t]
\centering
\small
\caption{Correlation of various training metrics with downstream benchmark performance (averaged over 5 benchmarks). $\psi$ and $\free$ substantially outperform loss as predictors of model quality.}
\label{tab:monitor}
\begin{tabular}{@{}lcccc@{}}
\toprule
\textbf{Metric} & \textbf{190M} & \textbf{1B} & \textbf{7B} & \textbf{Average} \\
\midrule
Training loss $\intE$ & 0.38 & 0.41 & 0.36 & 0.38 \\
Spectral entropy $\entropy$ & 0.72 & 0.78 & 0.81 & 0.77 \\
Free energy $\free$ & 0.83 & 0.87 & 0.89 & 0.86 \\
Order parameter $\psi$ & 0.90 & 0.93 & 0.94 & 0.92 \\
\bottomrule
\end{tabular}
\end{table}

The order parameter $\psi$ achieves $r > 0.92$ average correlation with downstream performance, nearly tripling the predictive power of loss alone ($r \approx 0.38$). Free energy $\free$ is the second-best predictor ($r \approx 0.86$), confirming that the thermodynamic perspective captures aspects of training quality invisible to loss curves.

\paragraph{Practical monitoring protocol.}
Based on these findings, we recommend a three-signal monitoring protocol:
\begin{enumerate}[leftmargin=*,itemsep=0pt]
    \item \textbf{$\psi$ (primary)}: if $\psi$ stops increasing while loss continues decreasing, the model is entering a glass state. Continuing training is thermodynamically unproductive.
    \item \textbf{$\sigma(t)$ (secondary)}: if entropy production rate drops to near zero, the system has reached thermal equilibrium at the current temperature. Either decrease LR (cool further) or switch data (change the potential landscape).
    \item \textbf{$\free$ (tertiary)}: free energy integrates loss and compression quality. A ``good'' training step decreases $\free$; a ``wasted'' step decreases loss but increases entropy (or vice versa), leaving $\free$ unchanged.
\end{enumerate}

%-----------------------------------------------------------------------------
\subsection{Q5: Optimal Schedule from First Principles}
\label{sec:optimal}
%-----------------------------------------------------------------------------

\paragraph{Minimum entropy production principle.}
The optimal thermodynamic process between two states minimizes total entropy production~\citep{liu_tegmark2025, okanohara2026}. For a schedule $\temp(t)$ taking the system from $(\intE_0, \entropy_0)$ to target $(\intE^*, \entropy^*)$ in time $T$, optimality requires constant entropy production rate:
\begin{equation}
    \sigma(t) = \text{const} \quad \Longleftrightarrow \quad \frac{d\free/dt}{\temp(t)} = \text{const}.
\end{equation}

\paragraph{Derivation of Gaussian decay.}
Using our fitted state equation (Eq.~\ref{eq:state}) and the relationship between $\free$, $\temp$, and $\entropy$, we solve the constrained variational problem:
\begin{equation}
    \min_{\temp(t)} \int_0^T \sigma(t)\,dt \quad \text{s.t.} \quad \intE(T) = \intE^*, \quad \int_0^T C(t)\,dt = C_{\mathrm{budget}}.
\end{equation}
The solution takes the form (derivation in Appendix~\ref{app:derivation}):
\begin{equation}
    \temp^*(t) = \begin{cases}
        \temp_{\max} & t \leq t_s \\[4pt]
        \temp_{\max} \cdot \exp\!\left(-\frac{(t - t_s)^2}{2\tau_{\mathrm{opt}}^2}\right) & t > t_s
    \end{cases}
    \label{eq:gaussian}
\end{equation}
where $t_s$ is the end of the stable phase and $\tau_{\mathrm{opt}} = (T - t_s)/\sqrt{2\ln(\temp_{\max}/\temp_{\min})}$. The Gaussian shape is neither linear (WSD) nor sinusoidal (cosine) but emerges naturally from the physics: it decays slowly at first (maintaining quasi-equilibrium as long as possible), accelerates through the transition region, and slows near the target (avoiding overshooting past the optimal basin).

\paragraph{Validation.}

\begin{table}[t]
\centering
\small
\caption{Final validation loss across schedules at matched compute. $\Delta$ computed relative to WSD-Linear baseline. Results averaged over 3 seeds; $\pm$ indicates standard error.}
\label{tab:schedules}
\begin{tabular}{@{}lcccc@{}}
\toprule
\textbf{Schedule} & \textbf{190M Loss} & $\Delta$ & \textbf{1B Loss} & $\Delta$ \\
\midrule
Cosine         & [tbd] & +1.4\% & [tbd] & +1.8\% \\
WSD-Linear     & [tbd] & baseline & [tbd] & baseline \\
WSD-Exponential & [tbd] & $-0.3\%$ & [tbd] & $-0.5\%$ \\
\textbf{Gaussian (ours)} & [tbd] & $\mathbf{-2.1\% \pm 0.2}$ & [tbd] & $\mathbf{-1.7\% \pm 0.3}$ \\
\bottomrule
\end{tabular}
\end{table}

The Gaussian schedule consistently outperforms all baselines. The improvement is larger at 190M (where the ideal gas approximation is more accurate) and remains significant at 1B. The thermodynamic efficiency $\eta_{\mathrm{thermo}}$ of the Gaussian schedule is [tbd], compared to [tbd] for WSD and [tbd] for cosine.

\paragraph{From theory to practice.}
The Gaussian schedule has three hyperparameters: peak LR $\temp_{\max}$ (from $\mu$P transfer~\citep{mup}), stable phase fraction $t_s/T$ (default 0.8, matching standard WSD), and Gaussian width $\tau_{\mathrm{opt}}$ (determined by the state equation). We release \ours as a drop-in replacement: \textbf{5 lines of code} substituting the decay function in any training loop. See Appendix~\ref{app:gaussian} for implementation.
